\documentclass[11pt]{article}
\def\activityid{F09-F-v2}\usepackage[letterpaper,margin=.65in,top=.60in,bottom=.65in]{geometry}
\usepackage[T1]{fontenc}
\usepackage[default]{sourcesanspro}
\usepackage{microtype,tikz,array,tabularx,fancyhdr,amsmath,amssymb}
\usepackage[hidelinks]{hyperref}
\usetikzlibrary{calc,arrows.meta}
\setlength{\parindent}{0pt}\setlength{\parskip}{7pt}
\setlength{\headheight}{0pt}\setlength{\footskip}{26pt}\setlength{\emergencystretch}{2em}
\pagestyle{fancy}\fancyhf{}\renewcommand{\headrulewidth}{0pt}\renewcommand{\footrulewidth}{.4pt}
\fancyfoot[L]{\scriptsize Bellingham Math Circle / Week 9 / \activityid}
\fancyfoot[R]{\scriptsize \thepage}
\definecolor{ink}{gray}{.12}\definecolor{soft}{gray}{.95}\color{ink}
\newcommand{\PageTitle}[2]{{\fontsize{10}{12}\selectfont\bfseries WEEK 9\quad BOUNCE LAB\hfill #1\par}\vspace{3pt}{\fontsize{26}{29}\selectfont\bfseries #2\par}\vspace{2pt}}
\newcommand{\NameLine}{{\small Name: \rule{2.65in}{.4pt}\hfill Date: \rule{1.35in}{.4pt}\par}\vspace{4pt}}
\newcommand{\Task}[2]{\par\vspace{3pt}{\large\bfseries #1}\par #2\par}
\newcommand{\Subhead}[1]{\par\vspace{3pt}{\large\bfseries #1}\par}
\newcommand{\RuleBox}[1]{\par\begingroup\setlength{\fboxsep}{8pt}\colorbox{soft}{\parbox{\dimexpr\linewidth-16pt\relax}{#1}}\endgroup\par}
\newcommand{\WriteBox}[2]{\par\textbf{#1}\par\vspace{2pt}\begin{tikzpicture}\draw[black!40,line width=.5pt] (0,0) rectangle (\dimexpr\linewidth-.5pt\relax,#2);\end{tikzpicture}\par}
\newcommand{\StudentFooter}{\vfill{\small Try it with objects. Explain by pointing, talking, or drawing.}\par}
\newcommand{\RookBoard}[3]{\begin{tikzpicture}[x=#3,y=#3]\draw[step=1,black!45] (0,0) grid (#1,#2);\draw[thick] (0,0) rectangle (#1,#2);\node[font=\Large] at ({#1-.5},.5){$\star$};\draw[-{Latex},thick] (0,-.35)--(#1,-.35);\draw[-{Latex},thick] ({#1+.35},#2)--({#1+.35},0);\end{tikzpicture}}
\newcommand{\Table}[3]{\begin{tikzpicture}[x=#3,y=#3]\draw[step=1,black!25] (0,0) grid (#1,#2);\draw[thick] (0,0) rectangle (#1,#2);\fill (0,0)circle(3pt);\draw[-{Latex},very thick] (0,0)--(.7,.7);\node[below] at ({#1/2},-.1){width #1};\node[rotate=90,above] at (-.2,{#2/2}){height #2};\end{tikzpicture}}
\newcommand{\GraphNodes}{\coordinate(A)at(0,0);\coordinate(B)at(3,0);\coordinate(C)at(1.5,2.6);\coordinate(D)at(1.5,.95);}
\newcommand{\Kfour}[1]{\begin{tikzpicture}[scale=#1]\GraphNodes\draw[very thick](A)--(B)--(C)--(A)--(D)--(B) (D)--(C);\foreach \n in {A,B,C,D}{\node[circle,draw,fill=white,minimum size=22pt,font=\bfseries]at(\n){\n};}\end{tikzpicture}}
\newcommand{\House}[1]{\begin{tikzpicture}[scale=#1]\coordinate(A)at(0,0);\coordinate(B)at(3,0);\coordinate(C)at(3,2);\coordinate(D)at(0,2);\coordinate(E)at(1.5,3.3);\draw[very thick](A)--(B)--(C)--(D)--(A) (D)--(E)--(C);\foreach \n in {A,B,C,D,E}{\node[circle,draw,fill=white,minimum size=22pt,font=\bfseries]at(\n){\n};}\end{tikzpicture}}

\begin{document}

\PageTitle{FACILITATOR / 1 OF 4}{From bouncing to straight lines}
\RuleBox{Core question: Predict the first corner and the number of bounces without tracing the whole path. Launch from bottom-left at 45 degrees. Stop at the FIRST corner; do not reflect there. Line crossings inside the table have no effect.}
\Subhead{Preparation, access, and the shared hour}
Print 3 K--1, 3 middle, 1 upper, and one extra reserve page. Bring seven counters, rulers, pencils, scrap grid paper, and two spare copies of the K--1 folding page. Folding is optional; an adult can reflect a traced room instead. Allow about 10 minutes to prepare; no actual ball is needed.

K--1: no reading or arithmetic; track a diagonal and choose one changed direction. Middle: count lengths and multiples to 12; coordinate two repeating lists. Upper: multiples, parity, and division; introduce gcd/lcm after a picture. The general coprime lemma additionally needs adult-supported subtraction of multiples and Euclid's argument. Extra: fractions and coordinate pairs; irrationality of $\sqrt2$ may be supplied as known, with its proof available on page 3.

0--10 explore counters and mirrored grid rooms. 10--15 brief reflection launch: which direction is blocked by a right wall? 15--35 main investigation: parent models K--1 while organizer launches middle and upper, then alternates visits; children trace, fold, and predict. 35--40 movement/reset. 40--55 continue with target designs or an optional proof. 55--60 share a prediction with its reason and tidy.

Triplets rotate tracer, predictor, and checker but each child should move a token. Parent mostly stays with K--1. The fifth grader receives a personal discussion, with a concrete next table left between visits.
\Subhead{Launch and hints}
``What if we opened the wall into a mirror room and kept going straight?'' Let the child fold a drawn path before introducing number theory. Ask which copied walls it reaches together. If tracing is difficult, adult draws while child chooses directions. Upper hint order: compare a rectangle with a scaled copy; list common multiples; count whole rooms; use odd/even; only then reduce by gcd.

For returners, ask for one remembered corner prediction, then go straight to mirrored rooms and explain it. Skip familiar traces and redundant table rows; the new destination is unfolding, an obstruction, or inverse design.

\Subhead{A satisfying stop; an optional proof}
\textbf{K--1:} predict and fold one route; continue by explaining the different 2-by-4 and 4-by-2 corners. \textbf{Middle:} build and explain a target table; continue by justifying the first common wall multiple. \textbf{Upper:} rule out bottom-left by an earlier corner and count walls; continue with adult-supported coprime division and the complete inverse design. Do not require every table row.

\newpage

\PageTitle{FACILITATOR / 2 OF 4}{Checked paths and corner rule}
\Subhead{K--1 and middle answers}
For width $w$, height $h$: $2\times2$ ends top-right with no bounce. $2\times4$ follows $(0,0),(2,2),(0,4)$, ending top-left; $4\times2$ follows $(0,0),(2,2),(4,0)$, ending bottom-right. Folding the two width-2, height-4 rooms turns the straight diagonal through $(2,2)$ into the first of these paths.

For $2\times3$, the folded path is $(0,0),(2,2),(1,3),(0,2),(2,0)$. In copied rooms, the first simultaneous wall meeting is $(6,6)$. The four middle rows have $(L,\text{widths},\text{heights},\text{corner})$:
\begin{center}\renewcommand{\arraystretch}{1.35}\begin{tabular}{c|c|c|c|c}Table&$L$&Widths&Heights&Corner\\\hline$2\times3$&6&3&2&bottom-right\\$3\times6$&6&2&1&top-left\\$3\times9$&9&3&1&top-right\\$4\times6$&12&3&2&bottom-right\end{tabular}\end{center}
A middle design reaching top-right after a bounce can be $2\times6$ (two bounces); top-left can be $2\times4$. Accept any verified design.
\Subhead{Upper theorem with all hypotheses}
Take positive integer $w,h$ and a 45-degree launch. Reflections unfold to $y=x$. A corner occurs when the common coordinate is both a multiple of $w$ and of $h$. Its first positive value is $L=\operatorname{lcm}(w,h)$. First use room counts $m=L/w,n=L/h$. If they shared $d>1$, distance $L/d$ would already cross $m/d$ widths and $n/d$ heights: an earlier corner. Thus they are coprime; both cannot be even, ruling out bottom-left without a gcd formula.

\textbf{Optional arithmetic bridge.} Write $w=ga,h=gb$ with coprime $a,b$. Any common multiple is $ga\,t$ with $b\mid at$. To justify $b\mid t$, begin with $a=2,b=3$: both $3t$ and $2t$ divide into threes, so their difference $t$ does too. In general, run repeated subtraction on strips of lengths $a,b$ and, alongside them, $at,bt$. Subtract the shorter from the longer until equal. Common divisors are unchanged, so the coprime base strips end at 1 and the scaled strips at $t$. Both scaled lengths remain multiples of $b$ throughout; hence $b\mid t$. The least positive $t$ is therefore $b$, and $L=gab$, $m=b,n=a$. The shared plan scaffold supplies physical bundle examples. If this step is supplied, record it as a given theorem.

An odd number of whole widths folds to the right, even to the left. An odd number of whole heights folds to the top, even to the bottom. Thus odd $a$, odd $b$ gives top-right; even $a$, odd $b$ gives bottom-right; odd $a$, even $b$ gives top-left. Both cannot be even, so first arrival at bottom-left is impossible. Scaling both dimensions preserves the reduced pair and corner.

Upper rows: $4\times6$: $L=12,m=3,n=2$, bottom-right; $6\times9$: $18,3,2$, bottom-right; $6\times10$: $30,5,3$, top-right; $8\times14$: $56,7,4$, bottom-right. Raw parity of the original lengths does not determine the endpoint.

\newpage

\PageTitle{FACILITATOR / 3 OF 4}{Bounces, inverse design, and slopes}
\Subhead{Count crossings, then design}
Before the final corner the line crosses $m-1$ vertical room walls and $n-1$ horizontal ones. With the optional lemma these are $b-1$ and $a-1$. No crossing belongs to both lists: that would be an earlier corner. So the number of bounces is $a+b-2$ (not counting launch or terminal corner). Path length, if wanted, is $\sqrt2 L$; $L$ itself is horizontal distance, not path length.

For top-right after exactly four bounces, $a,b$ are coprime odd positives with $a+b=6$. Candidates $(1,5),(3,3),(5,1)$ reduce to just $(1,5),(5,1)$ because $(3,3)$ is not coprime. Any positive integer scaling works, such as $2\times10$. Top-right after three bounces is impossible: two odd numbers have even sum, so $a+b-2$ is even. The latter is a parity obstruction, not a failed search.
\Subhead{Extra page: rational and irrational directions}
The extra uses a \emph{unit square} and physical slope $s=p/q$ in lowest terms. A positive lattice point $(u,v)$ on $y=(p/q)x$ satisfies $qv=pu$, so its first occurrence is $(q,p)$. By the coprime lemma on page 2, $q$ divides $u$; the first positive choice is $u=q$, giving $v=p$. The endpoint's horizontal side is right iff $q$ is odd; its vertical side is top iff $p$ is odd. There are $p+q-2$ bounces. Slope $2/3$ ends bottom-right with 3 bounces; $3/2$ ends top-left with 3.

For slope $\sqrt2$, an integer corner $(u,v)$ with $u>0$ would give $\sqrt2=v/u$, impossible. If requested, prove irrationality: reduce a supposed fraction $v/u$ first. From $v^2=2u^2$, $v$ is even; substituting $v=2k$ forces $u$ even, contradicting the reduction.

The other launch $(0,1/2)$ with slope 1 cycles through $(1/2,1),(1,1/2),(1/2,0),(0,1/2)$. It avoids corners but follows just a diamond. No-corner does not imply dense. Density of the $\sqrt2$-slope trajectory launched from $(0,0)$ is a separate theorem. Under our stopping rule, arbitrary starts need care: from $(0,2-\sqrt2)$ the same slope reaches unfolded $(1,2)$, folds to bottom-right, and stops.
\Subhead{Keep the conventions straight}
A general rectangle with physical slope $s$ becomes a square after rescaling axes, with normalized slope $sw/h$. Rationality of that normalized slope controls the square model. The extra fixes $w=h=1$ to avoid silently mixing those slopes. ``Rational polygon'' in research means its angles are rational multiples of $\pi$; it does not mean a rational launch slope. Our corner-start paths stop when they hit a corner, so do not call those stopped paths periodic billiard orbits.

\newpage

\PageTitle{FACILITATOR / 4 OF 4}{A doorway into dynamics}
\Subhead{Mathematical direction and primary sources}
Unfolding replaces a piecewise reflected path by a straight line on a lattice. The corner problem is divisibility, coprime integers, and parity; the inverse problem asks for a complete classification of parameters. These are substantive undergraduate number-theoretic and geometric ideas in a tangible representation.

Masur and Tabachnikov, \emph{Rational billiards and flat structures}, published in \emph{Handbook of Dynamical Systems} 1A (2002), pp. 1015--1089. In the linked author manuscript, \S1.3 gives unfolding, \S1.4 develops the unit-square/torus model, and \S1.5 extends the construction to flat surfaces. Later \S4 concerns periodic orbits. \href{https://math.uchicago.edu/~masur/handbook1.pdf}{Author manuscript}.

The classroom genuinely performs the first construction used in the research: reflect copies, then study straight motion. It does not establish the deeper recurrence or ergodicity theorems. The extra explicitly separates avoiding corners from density; the latter is supplied specifically for the origin-start $\sqrt2$ path under our stopping convention.
\Subhead{Downloaded sources and documented prior use}
JRMF \emph{Billiards Geometry}, PDF p. 2, asks for corner, bounce-count, and path-length rules. Lowell Handout 10, problems 10.2--10.7, includes $7\times5$, $4\times3$, $8\times10$, scaling, consecutive sides, $1\times n$, and $2\times\text{odd}$ families. Handout 11, problems 11.4--11.5, includes height-five rectangles and plus-shaped tables. Do not describe small rectangle families or scaling as wholly new for returning children.

Our changed emphasis is a proof through physical unfolding, an exact first-corner criterion, a bounce formula, and inverse existence/impossibility design. This deliberately promotes an old reserve into the core. Plus-shaped tables remain a later branch; do not claim the rectangle formula covers them.

\emph{Math Circle by the Bay}, preface printed pp. viii--x, supports manipulatives and revisiting ideas at different depths. Rozhkovskaya, Lesson 3 ``At the lesson,'' EPUB \texttt{part0013.xhtml}, supports attempts and child explanations before formal displays; Lesson 8's account warns about copying pace. Preprinted grids and optional adult tracing are our response, not a procedure claimed to be in those sources.
\Subhead{Keep a useful history}
IDs F09-K/M/U/X-v2. Record actual dimensions, launch slopes, whether unfolding was used, which formulas were explained, and which claims stayed conjectural. A prepared extension is not a taught activity. Future variants can change slopes, starting points, or polygons instead of enlarging the same arithmetic table.

\texttt{verify.py} independently traces all 900 rectangles with sides 1--30, compares the corner and bounce formulas, and checks the inverse-design list. The general proofs are on pages 2--3; finite checks alone do not establish them.

\end{document}
